% TOP_PROBLEM: {"id": 30006676, "title": "Generalized Riemann Hypothesis for Dirichlet L-functions", "category_id": 1, "category": {"id": 1, "name": "number_theory", "display_name": "Number Theory", "description": "Properties of integers, prime numbers, Diophantine equations.", "slug": "number-theory", "order_index": 1, "created_at": "2026-07-31T15:26:25.670Z"}, "status": "open"}
% FIRST_POSTED: 2026-09-27
% !TeX program = lualatex
% ATTEMPT_STATUS: unresolved
% Research continuation by GPT_6_Astra_Ultra, 27 September 2026.
% Catalogue statement and existing sources preserved verbatim.
\documentclass[11pt]{article}
\usepackage[margin=25mm]{geometry}
\usepackage{fontspec}
\setmainfont{Latin Modern Roman}
\usepackage{amsmath,amssymb,mathtools}
\usepackage{unicode-math}
\setmathfont{Latin Modern Math}
\usepackage{xurl,hyperref}
\hypersetup{colorlinks=true,urlcolor=blue}
\setcounter{secnumdepth}{0}
\setlength{\parindent}{0pt}
\setlength{\parskip}{5pt}
\setlength{\emergencystretch}{3em}
\begin{document}
\section*{\texorpdfstring{Generalized Riemann Hypothesis for Dirichlet L-functions}{Generalized Riemann Hypothesis for Dirichlet L-functions}}\label{30006676}
\subsection{Definitions and mathematical statement}
A Dirichlet character modulo $q$ is a homomorphism $({\mathbb{Z}}/q{\mathbb{Z}})^\times\to{\mathbb{C}}^\times$, extended periodically to integers and by zero on nonunits. Initially, for $\Re s>1$,
\[
 L(s,\chi)=\sum_{n\ge1}\frac{\chi(n)}{n^s}=\prod_p(1-\chi(p)p^{-s})^{-1}.
\]
Take the meromorphic continuation. For every $q\ge1$ and every character, ask whether
\[
 L(\rho,\chi)=0,\quad0<\Re\rho<1\quad\Longrightarrow\quad\Re\rho=\tfrac12.
\]
Principal and imprimitive characters remain in scope. The restriction to the open strip excludes poles and zeros on its boundary caused by removed Euler factors.
\par\smallskip\textit{Formulation and concept references:} \href{https://dlmf.nist.gov/25.15}{[S1]}.

\subsection{Short English statement}
Do all Dirichlet $L$-functions place their zeros inside the critical strip on its central vertical line?
\subsection{Research attempt}

\paragraph{Scope, literature check, and chosen attack.}
This attempt retains every modulus, including the principal character modulo
one. It therefore contains the ordinary Riemann Hypothesis as a special case.
The analytic continuation, primitive-character functional equation, and
imprimitive factorization used below are precisely the standard formulas
in [S1], equations 25.15.2, 25.15.4, and 25.15.5. An imprimitive character
induced by $\chi^*$ satisfies
\[
 L(s,\chi)=L(s,\chi^*)\prod_{p\mid q}(1-\chi^*(p)p^{-s}).
\]
Every nonconstant additional factor has its zeros on $\Re s=0$, since
$|\chi^*(p)|=1$. Thus reduction to primitive characters loses no zero in the
open strip. Principal characters reduce to $\zeta$ in the same way.

Reciprocal Dirichlet series provide a concrete arithmetic route to excluding
zeros. Garg--Maji [E1], Theorem 2.2, develop established smoothed
M\"obius-sum criteria for GRH. The present calculation uses a different,
unsmoothed weighted energy and proves only the sufficient implication it
needs. It does not claim a new equivalence theorem. The catalogue's further
references remain recorded, but no claimed resolution in [S3] is assumed.
The immediate objective is to turn the assertion into an explicit
cancellation estimate whose failure or proof can be investigated without
starting from a presumed list of critical-line zeros.

\paragraph{Reciprocal coefficients and a rigorous energy criterion.}
Write $\mu(n)=0$ when $n$ has a squared prime factor, and
$\mu(n)=(-1)^r$ when $n$ is a product of $r$ distinct primes; $\mu(1)=1$.
For a fixed character set
\[
 M_\chi(x)=\sum_{n\le x}\mu(n)\chi(n),\qquad
 \mathcal E_\chi(\sigma)=\int_1^\infty
       |M_\chi(x)|^2x^{-2\sigma-1}\,dx.
\]
The proposed missing estimate is
\begin{equation}\label{ra:9:energy}
 \mathcal E_\chi(\sigma)<\infty
 \quad\hbox{for every character $\chi$ and every $\sigma>1/2$.}
\end{equation}
Constants and the values of these integrals may depend on $q,\chi,\sigma$.
No bound uniform in the conductor is required by this particular criterion.

\textbf{Lemma 9.1.} If $\mathcal E_\chi(\sigma)<\infty$ for one real
$\sigma>0$, then $L(s,\chi)$ has no zero in $\Re s>\sigma$.

\textit{Proof.}
For $\Re s>1$, absolute convergence of the Euler product gives
\[
 \frac1{L(s,\chi)}=\sum_{n\ge1}\frac{\mu(n)\chi(n)}{n^s}
          =s\int_1^\infty M_\chi(x)x^{-s-1}\,dx.
\]
The last identity follows by summing
$s\int_n^\infty x^{-s-1}\,dx=n^{-s}$, with absolute convergence.
For $u=\Re s>\sigma$, Cauchy--Schwarz gives
\begin{equation}\label{ra:9:tail}
 \left|s\int_X^\infty M_\chi(x)x^{-s-1}\,dx\right|
 \le \frac{|s|X^{\sigma-u}}{\sqrt{2(u-\sigma)}}
 \left(\int_X^\infty |M_\chi(x)|^2x^{-2\sigma-1}\,dx\right)^{1/2}.
\end{equation}
Consequently the integral defines a holomorphic function $F_\chi$ on that
half-plane: on each compact subset the right side tends uniformly to zero,
so the truncated holomorphic integrals converge locally uniformly.
The meromorphic identity $F_\chi(s)L(s,\chi)=1$ holds first to the right
of both $1$ and $\sigma$, and then throughout the connected half-plane by
the identity theorem. A zero of $L$ there would make the product zero,
contradicting this identity. For a principal character, the pole at $1$
is harmless: the identity instead forces $F_\chi(1)=0$ when $1$ lies in
the domain. No division by $L$ near an unknown zero was used.
\hfill$\square$

If \eqref{ra:9:energy} holds, Lemma 9.1 excludes every strip zero with
real part greater than $1/2$. For a primitive character the functional
equation reflects a strip zero to one of the conjugate character across
$1/2$; its gamma factors are finite and nonzero inside the strip. The
family in \eqref{ra:9:energy} includes that conjugate character. Thus a
zero to the left would give a forbidden zero to the right. Together with
the imprimitive reduction, this proves that \eqref{ra:9:energy} suffices
for the full catalogue assertion. Conversely, an off-line zero with real
part $\beta>1/2$ forces $\mathcal E_\chi(\sigma)=\infty$ for every
$1/2<\sigma<\beta$. This contrapositive is a useful diagnostic, but no
converse from GRH to this energy condition is claimed here.

\paragraph{An arithmetic form of the missing cancellation.}
For a reduced residue $a$ modulo $q$ define the integer
\[
 M_{q,a}(x)=\sum_{\substack{n\le x\\n\equiv a\ (q)}}\mu(n),
 \qquad V_q(x)=\sum_{a\in(\mathbb Z/q\mathbb Z)^\times}M_{q,a}(x)^2.
\]
Character orthogonality, applied to the finite sums, gives exactly
\begin{equation}\label{ra:9:parseval}
 \sum_{\chi\bmod q}|M_\chi(x)|^2=\varphi(q)V_q(x).
\end{equation}
Indeed, after expanding the absolute square, the inner sum
$\sum_\chi\chi(m)\overline{\chi(n)}$ is $\varphi(q)$ if $m,n$ are units
with $m\equiv n\pmod q$, and zero otherwise. This also proves the identity
for complex characters; replacing an absolute square by an ordinary square
would not work.

A specific sufficient target is therefore
\begin{equation}\label{ra:9:variance}
 \forall q\ge1\ \forall\eta>0\ \exists C_{q,\eta}<\infty:
 \qquad V_q(x)\le C_{q,\eta}x^{1+\eta}\quad(x\ge1).
\end{equation}
For $\sigma=1/2+\delta$, take $\eta=\delta$. Each term on the left of
\eqref{ra:9:parseval} is nonnegative, so its integral is bounded by a
constant times $\int_1^\infty x^{-1-\delta}\,dx<\infty$. This proves
\eqref{ra:9:energy}. The estimate is required separately for every fixed
$q$; a statement for most moduli, or for moduli growing with $x$, would not
automatically supply it.

Expanding $V_q$ exposes the problem rather than resolving it:
\begin{align*}
 V_q(x)&=D_q(x)+C_q(x),\\
 D_q(x)&=\sum_{\substack{n\le x\\(n,q)=1}}\mu(n)^2\le x,\\
 C_q(x)&=2\sum_{\substack{1\le m<n\le x\\m\equiv n\ (q)\\(mn,q)=1}}
                    \mu(m)\mu(n).
\end{align*}
The diagonal already has the correct scale. The unresolved part is a
signed off-diagonal estimate $C_q(x)\ll_{q,\eta}x^{1+\eta}$.
Bounding every product by its absolute value gives only an $O_q(x^2)$
estimate. Positivity of $V_q$ supplies a lower bound $C_q\ge-D_q$,
not the needed upper bound. Thus the sum-of-squares reformulation does
not create the missing cancellation for free.

There is also no easy escape through averaging characters. If $\chi_0$
is principal, then
\[
 V_q(x)\ge\frac1{\varphi(q)}
       \left|\sum_{\substack{n\le x\\(n,q)=1}}\mu(n)\right|^2.
\]
This is Cauchy--Schwarz on the residue sums. At $q=1$ it is an equality
and the problem becomes the ordinary Mertens cancellation problem.
A large family average cannot justify deleting the principal character
or an exceptional member while retaining the stated target.

\paragraph{Stress tests of two tempting shortcuts.}
For a nonprincipal $\chi$, the untwisted sums $\sum_{n\le x}\chi(n)$
are bounded by $q$, because complete periods sum to zero. This proves
analytic continuation of the series for $L(s,\chi)$ into $\Re s>0$
by partial summation. It says nothing comparable about its reciprocal:
$\mu(n)\chi(n)$ is not periodic. For instance, for the real character
modulo four, $\mu(1)\chi(1)=1$ whereas $\mu(5)\chi(5)=-1$, despite
$1\equiv5\pmod4$. Confusing these two coefficient sequences replaces
zero-freeness by holomorphy, a strictly different assertion.

A second shortcut is to extrapolate an energy estimate from finite
partial sums. For any observed prefix of a bounded sequence, extend its
unseen tail first by alternating $+1,-1$, and alternatively by all $+1$.
The first extension has bounded partial sums and finite energy for every
$\sigma>0$; the second has partial sums asymptotic to $x$ and infinite
energy for $\sigma\le1$. The two sequences agree on the entire prefix.
These are auxiliary sequences, not alternative M\"obius functions, but
they show exactly why the omitted arithmetic tail needs a theorem.
Even testing only dyadic endpoint values is inadequate for general bounded
coefficients: on $[2^k,2^{k+1})$ put $+1$ on its first half and $-1$
on its second. Endpoint sums return to zero while interior sums grow
linearly with the block length. The integral records those excursions.

\paragraph{Boundary and uniformity audit.}
There is a useful concrete imprimitive check. Induce the character modulo
four to modulus twelve. Its $L$-function is
$L(s,\chi_4)(1+3^{-s})$, since the newly deleted Euler factor corresponds
to $\chi_4(3)=-1$. The extra zeros satisfy
$\Re s=0$ and $\Im s=(2k+1)\pi/\log3$, for integers $k$.
They do not contradict either Lemma 9.1 with $\sigma>1/2$ or the
catalogue's explicitly open-strip formulation. Trying to prove that the
entire imprimitive function has only critical-line zeros would instead
target a false statement. Similarly, allowing a pole of a principal
$L$-function at one does not invalidate a criterion for its reciprocal;
the reciprocal vanishes there, which is entirely compatible with
holomorphy.

The order of quantifiers deserves a separate check. Suppose for each
modulus an energy bound were available only for
$\sigma>1/2+1/q$. For that fixed modulus the lemma would leave an
uncontrolled strip of positive width. Sending $q$ to infinity does not
repair the missing statement for the original character. Conversely,
constants growing arbitrarily rapidly with $q$ would still suffice if
the exponent condition held for every positive $\delta$ and every fixed
$q$. This attempt needs neither a conductor average nor a single
constant controlling all characters. Finally, \eqref{ra:9:tail} is an
effective truncation bound only when its unknown tail energy is bounded
independently. Computing the finite integral and reusing its value as
a bound for the infinite tail would be circular.

\paragraph{Finite verification and reproducibility.}
The accompanying standard-library Python script computes $\mu$ by exact
prime factorization and checks \eqref{ra:9:parseval} at every integer
$1\le n\le4096$, for $q=1,3,4,5$. The four characters modulo five include
both genuinely complex characters, with values represented by Gaussian
integers. The step-function identity
\[
 \int_1^X V_q(x)x^{-2}\,dx
     =\sum_{n=1}^{X-1}\frac{V_q(n)}{n(n+1)}
 \qquad(X\in\mathbb Z,\ X\ge2)
\]
allows exact rational summation. Rounded display values are
\[
\begin{array}{c|rrrr}
q\backslash X&64&256&1024&4096\\\hline
1&1.317265&1.407570&1.455461&1.505754\\
3&2.400019&2.701077&2.870196&2.969291\\
4&2.223627&2.805885&3.110989&3.272757\\
5&3.023042&3.432138&3.681850&3.813208
\end{array}
\]
These are finite integrals at the endpoint exponent $\sigma=1/2$.
The argument requires infinite integrals at every strictly larger
exponent and asserts no endpoint convergence. The table is a check of
the bookkeeping, not evidence establishing convergence or a zero count.
The local record is
\texttt{\_Local/top\_problems/continuation\_2026\_09\_27/analytic/checks.py}
and its \texttt{results.json}; it contains no unreported zero search.

\paragraph{First gap and outcome.}
The analytic implication and the finite Fourier identity are proved here;
no simplicity, nonvanishing central value, or conductor-uniform bound is
assumed. The first missing assertion is \eqref{ra:9:variance}, or a
weaker estimate sufficient to make its weighted integral finite. Neither
the diagonal calculation nor periodic character cancellation bounds the
off-diagonal M\"obius correlations at the required scale. A useful next
attack would have to retain their signs, quantify the tail in
\eqref{ra:9:tail}, and cover each fixed modulus, including one. A disproof
along this route would still require a certified off-line zero; divergence
at $\sigma=1/2$ alone would not suffice.

\textbf{Outcome: UNRESOLVED.} This notebook gives a fully justified
sufficient energy criterion, its exact residue-class reduction, and
explicit failures of two proposed shortcuts. The needed arithmetic
cancellation has not been established, and no counterexample to GRH has
been produced. The process used direct derivation, primary-source checks,
and finite exact diagnostics; no external mathematical review is claimed.

\subsection{Sources}
\begin{itemize}
\item[S1]NIST DLMF25.15, Dirichlet L-functions, equations25.15.1–4 and zeros; GRH convention in Kandhil–Languasco–Moree, Mathematische Annalen394:43(2026).
\url{https://dlmf.nist.gov/25.15}.
\textit{Formulation source. Consulted 24 September 2026: Dirichlet characters, L-functions and zeros.}
\item[S2]Large values of quadratic character sums — Dong, Wang, Wang, Zhang and Zhao.
\url{https://arxiv.org/abs/2608.15773}.
\textit{Further reference; not a proof endorsement.}
\item[S3]On the Extended, Generalized, and Grand Riemann Hypotheses: A Unified Approach Based on the General Properties of L-Functions — Weicun Zhang, v19.
\url{https://www.preprints.org/manuscript/202506.0481}.
\textit{Further reference; not a proof endorsement.}

\item[E1]Meghali Garg and Bibekananda Maji,
\emph{Hardy--Littlewood--Riesz type equivalent criteria for the Generalized
Riemann hypothesis}, arXiv:2208.07596v1 (16 August 2022), Theorem 2.2.
\url{https://arxiv.org/pdf/2208.07596v1}.
\textit{Primary research context for smoothed reciprocal-coefficient criteria.
Consulted 27 September 2026; not used as a premise of Lemma 9.1.}
\end{itemize}
\end{document}
